The Adams method represents the small-state extreme of the divisor
method spectrum.
Its ceiling rounding guarantees that every state receives at least
its upper quota entitlement, and its maximisation of minimum
per-capita representation makes it the most protective method
for small states facing the rounding problem.

The paradox immunity proof (Theorem~\ref{thm:adams-paradox})
establishes that even the most extreme small-state bias does not
compromise the monotone allocation property that makes divisor methods
superior to Hamilton.
This is an important theoretical point: paradox immunity is a consequence
of the divisor method structure, not of any particular rounding rule.
Adams, HH, Webster, and Jefferson are all paradox-immune despite
spanning the entire range of small-to-large-state bias.

The 2020 Census analysis shows that Adams and HH currently agree,
reflecting the fortuitous circumstance that no state has an exact quota
in the (narrow) range where the methods differ.
As populations shift in future censuses, this agreement may not persist.
In smaller historical house sizes ($H < 400$) or future censuses with
different population distributions, Adams and HH may diverge: with fewer
seats the exact quotas shift, and small states with quotas in the range
$(s, \sqrt{s(s+1)})$ would gain a seat under Adams but not under HH.
\textsc{bisect-apportion}'s Adams implementation enables practitioners
to monitor whether Adams and HH would diverge for future census data
and to compute the counterfactual apportionment for advocacy,
policy analysis, and litigation support.

The theoretical significance of Adams extends beyond its direct applicability:
it defines one end of the fairness spectrum that Congress implicitly
navigated when adopting Huntington-Hill in 1941.
Understanding where HH sits on the Adams-to-Jefferson spectrum ---
closer to Adams than to Webster --- is essential context for
evaluating the Supreme Court's \textit{Montana} holding that
Congress's choice was rational.
